\subsubsection{Итерация [[K]]}

\begin{enumerate}
  \item Найдём значение градиента $\omega^{([[K]])}=\textrm{grad}\,f(x^{([[K_PREV]])})$:
  \[
  \begin{gathered}
    \omega^{([[K]])}=(2ax_1^{([[K_PREV]])}+cx_2^{([[K_PREV]])}+d,2bx_2^{([[K_PREV]])}+cx_1^{([[K_PREV]])}+e)^T=\\
    =(2\cdot [[A]]\cdot [[X_PREV]]+[[C]]\cdot [[Y_PREV]]+[[D]], 2\cdot [[B]]\cdot [[Y_PREV]]+[[C]]\cdot [[X_PREV]]+[[E]])^T=\\
    =[[OMEGA]]
  \end{gathered}
  \]

  \item Вычислим \( x^{([[K]])} \):
  \[
    x^{([[K]])}=x^{([[K_PREV]])}-\eta\omega^{([[K]])}=([[X_NEW]],[[Y_NEW]])^T
  \]

  \item Графическая интерпретация:
  \begin{figure}[H]
  \centering
  \begin{tikzpicture}
    \begin{axis}[
      width=10cm, height=10cm,
      axis lines=middle,
      xlabel={$x$},
      ylabel={$y$},
      xlabel style={right},
      ylabel style={above},
      colormap={blackwhite}{gray(0cm)=(0.4); gray(1cm)=(0.9)},
      view={0}{90}
    ]
      \addplot3[
        contour lua={
          levels={[[Z_LEVELS]]},
          labels=false,
        },
        samples=100,
        domain=[[X_MIN]]:[[X_MAX]],
        domain y=[[Y_MIN]]:[[Y_MAX]],
      ] {[[A]]*x^2 + [[B]]*y^2 + [[C]]*x*y + [[D]]*x + [[E]]*y + [[F]]};

      [[GD_PATH_ITER]]
        
    \end{axis}
  \end{tikzpicture}
  \caption{Итерация №[[K]] метода градиентного спуска.}
  \end{figure}

  \item Проверка критерия останова:
  \[
    [[BREAKPOINT_CHECK]]
  \]
\end{enumerate}
